\section{Methods}
\label{sec:methods}

\subsection{Cohort}
\label{sec:methods-cohort}

The Kineret dataset spans six Israeli general hospitals, harmonized to a
common data model by the Kineret project. Extraction, cohort construction
and abstraction ran inside a regulated hospital computing environment; the
present study uses the same admission set as the concurrent adherence
study \citep{oded2027kineret-jbi}.

Eligibility. An admission was eligible if it was an inpatient visit with
at least one measurement recorded overnight (18:00--07:00, taken as
evidence of a genuine overnight stay); length of stay between $48$ hours
and $14$ days; patient at least $18$ years old at admission; the admission
did not occur in an excluded department (pediatric or obstetric); the
visit identifier was not associated with more than one patient identifier
(treated as a data-integrity failure); the patient did not die within
$48$ hours of admission; and the admission met at least one criterion for
diabetes or hyperglycemia. Diabetes or hyperglycemia was established by
at least one of: a documented diabetes diagnosis (from a curated concept
list); recurrent hyperglycemia (two glucose measurements $\ge 180$\,mg/dL
separated by $2$--$24$ hours); or extreme hyperglycemia (any glucose
$\ge 250$\,mg/dL). Hypoglycemia is not an entry criterion; it is treated
as an outcome throughout.

Restricting the eligible set to the study period of August $2022$ through
December $2025$ yields the analytic cohort of $48{,}354$ admissions among
$29{,}577$ patients. Full cohort characteristics are given in
Table~\ref{tab:cohort} of the concurrent JBI paper
\citep{oded2027kineret-jbi} and reproduced in Appendix~\TODO{A once
cross-references are frozen}; salient for prediction, the cohort is
older (mean age $71.9$\,y, median $73.3$\,y), predominantly non-ICU, with
a median length of stay of $5.55$ days and $30$-day mortality of
$9.58\%$.

\subsection{Targets}
\label{sec:methods-targets}

Ten complications are the candidate targets: $30$-day mortality,
hyperglycemia, severe hyperglycemia, hypoglycemia, severe hypoglycemia,
hyperosmolality, ketoacidosis, infection, inpatient acute kidney injury /
kidney complication, and inpatient myocardial infarction / cardiovascular
complication. Nine are emitted as event abstractions by the same Mediator
knowledge base that scores adherence; $30$-day mortality is taken from
linked death records. Operational definitions match the adherence study
\citep{oded2027kineret-jbi}; using one endpoint set across both analyses
lets a recommendation's measured association with an outcome and a
model's predictive use of that recommendation be read against one
another.

Two targets carry the labels the raw source table does not: the Mediator
output records complications as knowledge-base verdicts (e.g.,
\texttt{HYPERGLYCEMIA\_EVENT}), whereas the raw table records the
measurements those verdicts read. To bring the two into agreement, our
pipeline re-derives the raw side's own complication rows by applying, in
Python, the same threshold rules the Mediator's XML applies (with
patient-baseline ratios for AKI and a minimum-count guard for the
glycemic events), and then intersects the two views at $1$-minute
tolerance to keep only occurrences the two files date identically. The
result is a canonical event table used for labels and re-injected into
every arm's input stream, so event support is identical across arms by
construction. Cross-file rule agreement was audited manually against the
knowledge-base XML.

\subsection{The prediction task}
\label{sec:methods-task}

Each admission is observed over its first $K$ days and each of the eight
target complications that clear a prevalence threshold is predicted with
three heads: whether it occurs in the label window $(K \cdot 24,\,14 \cdot
24]$\,hours after admission (multi-label BCE risk), when it occurs (per-
outcome onset regression trained on positives only, in $z$-normalized
hours), and length of stay (regression trained on discharged admissions).
Both onset and length of stay are trained jointly with risk in every arm,
including logistic regression---this is a wider contract than the
MIMIC-IV study, where onset was an INTERVenE-only auxiliary.

$K$ enters twice, and the distinction matters. In \emph{training} it is a
data augmentation: each admission contributes one sample per context
window it reaches over $K \in \{1,\ldots,13\}$ days---the same trajectory
cut at a different point with its label window re-derived to match, so
the model learns to forecast from however much history it happens to
have. A patient discharged on day $N$ contributes $\min(13, \lceil N
\rceil - 1)$ samples---exactly $N - 1$ for whole-day stays---so long
stays contribute more samples than short ones, which is the correct
weighting since they contain more forecastable history.

In \emph{evaluation} $K$ is a single fixed window, chosen from the
target-support table before training begins, and every reported number
comes from it. Reporting at one window rather than sweeping keeps the
results a table rather than a surface, and keeps the head layout
fixed---targets are support-filtered once, on train, at that window. On
this cohort the chosen evaluation window is $K = 4$ days (justification
in Section~\ref{sec:results-support}), the same order of magnitude as
INTERVenE's MIMIC-IV window ($K = 2$ days) but longer, reflecting the
longer inpatient stays of the general-ward cohort.

Because the horizon end is fixed at $14$ days, $K$ also sets how far
ahead the model is forecasting: $K{=}1$ predicts over $13$ days, $K{=}13$
over one. Both numbers---current context length and remaining horizon---
are handed to the model as explicit context features, so this is not
leakage: they are constant at evaluation and known at prediction time by
construction.

\subsection{Cohort membership and splits}

A patient is in the study if they can furnish at least one training
sample---i.e., their stay outlasts the shortest training window. The
validation and test splits then take only the patients who reach
$K_{\mathrm{eval}} = 4$ days, because a $3$-day admission cannot be
scored at $K{=}4$. So a $3$-day stay trains the model (via its $K{=}1$
and $K{=}2$ cuts) without ever being scored, and raising the evaluation
window shrinks validation and test while leaving training untouched.
The alternative---keying membership on the evaluation window---would let
a reporting decision quietly evict short admissions from training whose
early cuts are perfectly good supervision. Cohort membership is decided
against the \emph{minimum} trajectory length across the raw and
Mediator-derived source tables rather than their union, since a union
would admit patients the raw-data arms could score but INTERVenE would
silently drop (the Mediator's last interval routinely closes before the
last raw measurement).

The split is over \emph{people}, not admissions
(\texttt{split\_group\_by\_person}): two admissions of one patient share
comorbidities, baseline physiology and often the same recurring
complication, so training on one and scoring the other measures
memorization of that patient rather than generalization to a new one.
The split is frozen at build time and every augmented sample of an
admission inherits its admission's split, so a sample can never cross
splits or appear on both sides of the augmentation.

Admissions whose label window would run past the end of the extract are
dropped, since scoring their unobservable late complications as
negatives is a censoring artefact indistinguishable from a true negative
(\texttt{date\_range\_require\_full\_horizon}).

\subsection{The seven-arm ladder}
\label{sec:methods-ladder}

Every arm predicts the same targets on the same admissions at the same
evaluation window with the same head layout; the arms differ only in
input representation. Table~\ref{tab:ladder-arms} lists them; the
contrasts they enable follow.

\begin{table}[htbp]
\centering
\small
\caption{The seven arms. Every arm shares the cohort, the labels, the
patient split, and the head layout. Only the input representation changes.}
\label{tab:ladder-arms}
\begin{tabularx}{\textwidth}{l>{\raggedright\arraybackslash}p{0.28\textwidth}>{\raggedright\arraybackslash}X}
\toprule
Rung & Model & Input \\
\midrule
1 & LogReg                      & Raw measurements, distribution-summarized; no temporal model \\
2 & LogReg $+$ QA               & As above, plus aggregated compliance in the context vector \\
3 & ss-STraTS                   & Raw measurements with learned temporal attention \\
4 & ss-STraTS $+$ QA            & As above, plus aggregated compliance in the context vector \\
5 & INTERVenE-Enc ($\sigma$-bins) & Knowledge-free intervals over the same raw data \\
6 & INTERVenE-Enc (KB)          & Knowledge-based Mediator abstractions \\
7 & INTERVenE-Enc (KB $+$ QA)   & As above, plus treatment-quality pattern tokens \\
\bottomrule
\end{tabularx}
\end{table}

Three contrasts carry the argument. Rungs $5 \to 6$ isolate the knowledge
base: same raw data, same architecture, same labels, different
abstraction layer. Rungs $6 \to 7$ isolate the treatment-quality signal.
Rungs $1 \to 2$ and $3 \to 4$ put the same QA question to models that
cannot represent pattern tokens, so a QA effect that appears only for
INTERVenE means something different from one that appears everywhere.

\subsubsection*{Arm 1: LogReg}
For each measurement concept, the observed values over $[0, K \cdot 24]$
hours are summarized by mean, standard deviation, minimum, maximum,
first, last, and count; missingness is indicated by a per-concept
observed-flag. The resulting fixed-length vector is concatenated with the
static context and passed through a multi-label logistic-regression
head. Onset and length-of-stay heads share the same input. This arm
discards temporal ordering entirely, so it is the floor: any arm that
does not clear it is not using time.

\subsubsection*{Arm 3: ss-STraTS}
The self-supervised STraTS \citep{tipirneni2022strats} baseline is used
with the pretraining stage removed. The MIMIC-IV benchmark showed
pretraining offered no consistent lift on this cohort and the two-stage
variant is not run inside the regulated environment. Every observation
is represented as a (concept, time, value) triplet; the sequence is
processed by set-based Transformer attention. Labels come from the
shared cohort; onset and length-of-stay heads are added, both trained
in the same schedule as the risk head. The head layout is $2K + 1$
outputs ($K$ risk logits, $K$ $z$-scored onset times, one length-of-stay)
so that timing is scored comparably across arms.

\subsubsection*{Arm 5: INTERVenE-Enc ($\sigma$-bins)}
The same INTERVenE-Enc architecture, but with the knowledge base
replaced by a knowledge-free discretization: each measurement is
bucketed into one of seven bands ($\pm 0.5\sigma$, $\pm 1\sigma$,
$\pm 2\sigma$; centred on the global training-set mean and standard
deviation for that concept) and consecutive same-concept observations
within $24$ hours are merged into a single interval so the encoder has
interval structure to model at all. Concepts are named
\texttt{<CONCEPT>\_STD\_<BIN>} and receive an augmented TAK repository
so INTERVenE's hierarchical embedder still resolves the concept
ancestry. This arm receives no knowledge-base-derived complication
events on its input, since without a knowledge base one would not have
them (\texttt{std\_include\_kb\_events} $= \mathrm{False}$).

\subsubsection*{Arms 6 and 7: INTERVenE-Enc (KB and KB $+$ QA)}
Arm~6 consumes the Mediator's output directly---the knowledge-based
temporal abstractions (states, trends, contexts, events) as
interval-typed tokens. Arm~7 adds the quality-assessment pattern tokens:
each formalized recommendation from the adherence study
\citep{oded2027kineret-jbi} emits an interval token
(\texttt{<PATTERN>\_START} / \texttt{<PATTERN>\_END}) wherever it applies,
carrying its compliance score in the value channel. Both arms
additionally receive the aggregated QA block in their static context
vector (see below).

\subsection{The quality-assessment signal is deliberately asymmetric}

\begin{table}[htbp]
\centering
\small
\caption{The QA signal is supplied in the richest form each architecture
admits. \emph{QA aggregate}: mean compliance per pattern over each
sample's $[0, K \cdot 24]$ window, appended to the static context vector,
identical across arms. \emph{Pattern tokens}: the compliance-scored
intervals themselves, available only to interval-consuming architectures.}
\label{tab:qa-asymmetry}
\begin{tabularx}{0.85\textwidth}{Xcc}
\toprule
Arm & QA aggregate block in context & \texttt{*\_PATTERN} interval tokens \\
\midrule
LogReg (raw + QA)                 & \checkmark & --- \\
ss-STraTS (raw + QA)              & \checkmark & --- \\
INTERVenE-Enc (KB + QA)           & \checkmark & \checkmark \\
\bottomrule
\end{tabularx}
\end{table}

Every QA arm receives the same aggregate block. Only INTERVenE additionally
receives the compliance patterns as temporal tokens, because it is the only
model that can represent them (they are intervals). Each architecture gets
the treatment signal in the richest form it admits: anything else would
either starve INTERVenE or credit the others with a representation they
cannot use. The aggregation is a mean of the pattern's compliance scores
over that sample's own $[0, K \cdot 24]$ window, so a $K{=}2$ sample never
sees four days of compliance and no arm reads compliance from the label
window.

\subsection{Knowledge containment and leakage guards}
\label{sec:methods-guards}

The Mediator emits \texttt{*\_EVENT} concepts that are not prediction
targets---for example \texttt{AKI\_EVENT} is a knowledge-base verdict
reached by reading raw creatinine---and handing one to LogReg or
ss-STraTS gives them a derived clinical judgement they could not have
computed themselves. Rungs $5 \to 6$, which are meant to isolate the
value of the knowledge base, would then measure a leak rather than
abstraction. The two streams therefore differ by design
(\texttt{kb\_events\_for\_kb\_arms\_only}):
target events are given to every arm as observed history, whereas
non-target \texttt{*\_EVENT} concepts are stripped from the raw stream
and kept in the KB stream. A unit test asserts both directions, and the
KB abstraction stream is harmonized with the raw stream so both see the
same reconciled canonical event support.

Terminal markers (\texttt{RELEASE}, \texttt{DEATH}) are stripped from
every input stream in every spelling. Outcome events falling inside the
observation window are not stripped: they are observed history, and the
label window starts strictly after the input window, so an event an arm
can see is history, never the thing it is being asked to predict.
Measurements from which outcomes are derived (glucose, creatinine,
etc.) are never stripped either way: they are legitimate predictors,
and dropping them would penalize exactly the arms that read raw data.

\subsection{Training}

Models are trained from scratch under a three-phase schedule
\citep{oded2027intervene}, not fine-tuned from the MIMIC-IV checkpoint.
The two cohorts differ in target set, patient distribution, and
knowledge base; a fine-tuned model would confound the architecture's
contribution with the initialization's, and training from scratch keeps
the seven arms comparable to one another.

The three neural architectures (LogReg, ss-STraTS, INTERVenE-Enc) are
all PyTorch models---the logistic-regression arm is a linear head with
the same optimizer, so all three arms run on the same device.
Per-model hyperparameters are those released in each model's reference
repository and are not tuned to this cohort; the per-cohort settings
are those governing the shared inputs (target list, evaluation window,
QA aggregation, event reconciliation), which live in a single
\texttt{configure\_study(\ldots)} call whose exact value is
\texttt{study\_config.json} in the reproducibility bundle.

\subsection{Evaluation}
\label{sec:methods-eval}

The primary metric is support-weighted AUPRC across the retained
targets. AUROC, best-$F_1$, best $F_1$ at $\tau{=}0.5$, minimum
recall-at-precision, onset MAE (hours), and length-of-stay MAE (hours)
are secondary. Per-outcome metrics are additionally reported for each
target. Uncertainty is a $2{,}000$-resample patient-level bootstrap; two
arms are treated as separated by the data only when their $95\%$
intervals do not overlap. This is sampling variability for the model
that was trained, not the variance of the training procedure, since
multi-seed retraining was not affordable at this scale (see
Section~\ref{sec:discussion-limits}).

Onset MAE is conditional on occurrence---only positive patients
contribute. For an outcome with a handful of positives the interval
reflects which of those few patients got drawn rather than model
uncertainty, and at $n = 1$ collapses to a point and reads as
precision, which is the opposite. Rows with fewer than five positives
per bootstrap draw are flagged \texttt{ci\_reliable = False}: the
estimates are kept but not to be quoted alone.

A \emph{wiring check} runs all seven arms for one epoch before any full
training, and asserts that every arm was scored on identical labels and
identical test patients. It fails loudly if not.
